\section{数据精炼}\label{sec:refine}

\begin{chapterquestion}
\textbf{本章研究的问题}：高不确定区域是否对应代理模型的真实认知盲区？
\end{chapterquestion}

\subsection{方法定位：不确定性引导的局部重采样}

标准主动学习（Active Learning）需要获取\emph{新的真实标签}（如新 CFD 仿真）。
本文受限于公开数据集，无法引入新物理标签，
因此本模块在严格意义上是\textbf{不确定性引导的局部重采样}
（Uncertainty-Guided Local Resampling）：
用 $\sigma(x)$ 识别代理薄弱区，从\emph{已有}训练集中按相似性检索补充样本后重训练。
本章的目的并非仅报告 MSE 的下降幅度，
而是借此检验一个判别性问题：\textbf{代理模型是否存在真实的认知盲区？}

\subsection{单轮反馈精炼流程}

在第~\ref{sec:risk}~节风险敏感 PSO 基础上，执行一次完整的反馈循环：
（1）PSO 探索并收集候选位置的 $\sigma(x)$；
（2）选取 Top-$20$ 个 $\sigma$ 最高的候选，即代理预测不确定性最高的位置；
（3）在这些位置附近从训练集检索最近邻样本（共 $100$ 个）加入训练；
（4）重训练 DeepEnsemble。

\begin{table}[H]
\centering
\caption{Top-5 高不确定性候选解的特征}
\begin{tabular}{rrrrrrr}
\toprule
速度 & 翼角 & DRS & 下压力 (N) & 阻力 (N) & $\mu$ & $\sigma$ \\
\midrule
175 & 5.1 & 0 & 8773 & 154.8 & 36.34 & 37.57 \\
217 & 14.7 & 1 & 6081 & 260.0 & 46.02 & 36.73 \\
139 & 34.6 & 0 & 6911 & 196.6 & 41.63 & 36.63 \\
106 & 40.4 & 0 & 6331 & 203.4 & 36.17 & 32.13 \\
114 & 22.3 & 0 & 9093 & 155.3 & 27.68 & 31.36 \\
\bottomrule
\end{tabular}
\label{tab:candidates}
\end{table}

表~\ref{tab:candidates} 显示高不确定候选解的共同特征：\textbf{下压力极高（6000--9000 N）}，
对应训练数据中极少出现的极端工况。这与第~\ref{sec:risk}~节 OOD 测试的结论一致——
高不确定区域对应数据覆盖不足之处。

\subsection{结果：高不确定区域的误差显著下降}

\begin{table}[H]
\centering
\caption{精炼前后模型性能}
\begin{tabular}{lrrr}
\toprule
指标 & 精炼前 & 精炼后 & 变化 \\
\midrule
R$^2$  & 0.9995 & 0.9998 & $+0.03\%$ \\
MSE   & 0.3272 & 0.1444 & $\mathbf{-55.9\%}$ \\
MAE   & 0.2712 & 0.1678 & $-38.1\%$ \\
RMSE  & 0.5720 & 0.3800 & $-33.6\%$ \\
\bottomrule
\end{tabular}
\label{tab:refine}
\end{table}

\begin{figure}[H]
\centering
\includegraphics[width=0.78\textwidth]{figures/active_learning_convergence.png}
\caption{精炼前后 PSO 收敛曲线对比}
\label{fig:al_conv}
\end{figure}

一轮精炼后，MSE 下降 \textbf{55.9\%}，PSO 收敛速度提升约 \textbf{2.1 倍}（$37\to18$ 代），
且收敛后的 gbest 从 $101.30$（含过度外推）降至 $100.02$（图~\ref{fig:al_conv}）。
误差下降主要发生在高 $\sigma$ 区域，gbest 的回落对应过度外推被纠正。
两者共同支持以下判断：

\begin{chapterquestion}
代理模型在低覆盖、高不确定区域存在认知盲区；
该盲区并非统计噪声，而是可被针对性修复的缺陷。
\end{chapterquestion}

\subsection{方法的局限}

需如实指出本方法的三点局限，以免过度解读：
\begin{enumerate}
    \item \textbf{隐式重加权}：检索到的最近邻若已存在于训练集，
        其效果等价于对这些样本施加更高训练频率（Instance Reweighting），
        而非注入新的分布信息；故 55.9\% 的 MSE 降幅应理解为
        OOD 区域\emph{拟合质量}的改善，不应等同于泛化能力的等量提升。
    \item \textbf{无消融对照}：本实验未设随机重采样对照组，
        因此不能完全排除"任何形式的样本扩充都会改善精度"的可能；
        严格的因果归因需消融框架完成。
    \item \textbf{仅单轮}：本文只执行了一轮"探索$\to$增强$\to$重训练"反馈，
        尚未构成多轮迭代的闭环；多轮收敛性与边际收益递减留待未来工作。
\end{enumerate}

至此，可信性分析的三个环节——度量（第~\ref{sec:uncertainty}~节）、
约束（第~\ref{sec:risk}~节）、修补（本节）——均已完成。
第~\ref{sec:converge}~节进一步考察异常二与异常三是否仍然存在。
